\subsubsection*{Invalid Argument Forms}

\begin{remark*}[Working vocabulary]
This topic uses counterexample assignments to show why familiar invalid forms
fail to preserve truth from premises to conclusion.
\end{remark*}

% BEGIN GENERATED ARTIFACT: thm:affirming-consequent-invalid-propositional-logic
\begin{theorembox}{Theorem (Affirming the Consequent is Invalid)}
\begin{theorem}[Affirming the Consequent is Invalid]
\label{thm:affirming-consequent-invalid-propositional-logic}
The argument form
\[
\varphi\to\psi,\qquad \psi \quad\therefore\quad \varphi
\]
is invalid.

\noindent\hyperref[prf:affirming-consequent-invalid-propositional-logic]{Go to proof.}
\end{theorem}
\end{theorembox}

\begin{remark*}[Standard quantified statement]
\[
\neg\operatorname{Valid}(\{\varphi\to\psi,\psi\},\varphi).
\]
\end{remark*}

\begin{remark*}[Predicate reading]
There is a truth assignment under which \(\varphi\to\psi\) and \(\psi\) are true
while \(\varphi\) is false.
\end{remark*}

\begin{remark*}[Interpretation]
A conditional can be true because both antecedent and consequent are true, but
also because the antecedent is false. Therefore truth of the consequent does not
force truth of the antecedent.
\end{remark*}

\begin{dependencies}
\begin{itemize}
  \item \hyperref[def:counterexample-assignment-propositional-logic]{Counterexample Assignment}
  \item \hyperref[def:valid-argument-form-propositional-logic]{Valid Argument Form}
\end{itemize}
\end{dependencies}
% END GENERATED ARTIFACT: thm:affirming-consequent-invalid-propositional-logic

% BEGIN GENERATED ARTIFACT: thm:denying-antecedent-invalid-propositional-logic
\begin{theorembox}{Theorem (Denying the Antecedent is Invalid)}
\begin{theorem}[Denying the Antecedent is Invalid]
\label{thm:denying-antecedent-invalid-propositional-logic}
The argument form
\[
\varphi\to\psi,\qquad \neg\varphi \quad\therefore\quad \neg\psi
\]
is invalid.

\noindent\hyperref[prf:denying-antecedent-invalid-propositional-logic]{Go to proof.}
\end{theorem}
\end{theorembox}

\begin{remark*}[Standard quantified statement]
\[
\neg\operatorname{Valid}(\{\varphi\to\psi,\neg\varphi\},\neg\psi).
\]
\end{remark*}

\begin{remark*}[Predicate reading]
There is a truth assignment under which \(\varphi\to\psi\) and \(\neg\varphi\)
are true while \(\neg\psi\) is false.
\end{remark*}

\begin{remark*}[Interpretation]
A conditional does not say that the antecedent is the only way for the consequent
to be true. The antecedent can fail while the consequent still holds.
\end{remark*}

\begin{remark*}[Exposition]
For both invalid forms, take a truth assignment with \(\varphi\) false and
\(\psi\) true. Then \(\varphi\to\psi\) is true. This assignment makes the
premises true and the conclusion false in each corresponding invalid pattern.
\end{remark*}

\begin{dependencies}
\begin{itemize}
  \item \hyperref[def:counterexample-assignment-propositional-logic]{Counterexample Assignment}
  \item \hyperref[def:valid-argument-form-propositional-logic]{Valid Argument Form}
\end{itemize}
\end{dependencies}
% END GENERATED ARTIFACT: thm:denying-antecedent-invalid-propositional-logic
